\documentclass[11pt]{article}
\usepackage[a4paper,margin=1in]{geometry}
\usepackage{fontspec}
\usepackage[boldfont=false]{xeCJK}
\setCJKmainfont{FandolSong-Regular.otf}
\usepackage{amsmath}
\usepackage{amssymb}
\usepackage{array}
\usepackage{booktabs}
\usepackage{graphicx}
\usepackage{adjustbox}
\usepackage{enumitem}
\setlist[itemize]{topsep=2pt,partopsep=0pt,parsep=0pt,itemsep=2pt,leftmargin=1.4em}
\setlist[enumerate]{topsep=2pt,partopsep=0pt,parsep=0pt,itemsep=2pt,leftmargin=1.6em}
\usepackage{parskip}
\usepackage{fvextra}
\fvset{breaklines=true,breakanywhere=true,fontsize=\small}
\setlength{\parindent}{0pt}

\begin{document}

\begin{center}
  {\LARGE\bfseries Movement into and out of cells}\\[0.35em]
  {\large Igcse Biology}
\end{center}
\vspace{0.6em}

\section*{Diffusion}

\textbf{Diffusion} 扩散 is the \textbf{net movement} 净移动 of \textbf{particles} 粒子 from a region of their higher \textbf{concentration} 浓度 to a region of their lower concentration. We say the particles move \textbf{down a concentration gradient} 浓度梯度.

The particles spread out because of their own \textbf{random movement} 随机运动. The \textbf{energy} 能量 for diffusion comes from the \textbf{kinetic energy} 动能 of the random movement of \textbf{molecules} 分子 and \textbf{ions} 离子 — so diffusion needs \textbf{no} extra energy.

Some substances pass into and out of cells by diffusion through the \textbf{cell membrane} 细胞膜.

\textbf{Why diffusion matters.} Many \textbf{gases} 气体 and dissolved substances (\textbf{solutes} 溶质) move by diffusion in living things. For example, \textbf{oxygen} 氧气 diffuses into cells for respiration, and \textbf{carbon dioxide} 二氧化碳 diffuses out.

The rate of diffusion changes with four factors:

\begin{center}
\adjustbox{max width=\linewidth}{%
\begin{tabular}{ll}
\toprule
\textbf{Factor} & \textbf{Diffusion is faster when…} \\
\midrule
\textbf{surface area} 表面积 & the surface is larger \\
\textbf{temperature} 温度 & it is hotter (particles move faster) \\
concentration gradient & the difference in concentration is bigger \\
distance & the distance to travel is shorter (a thinner barrier) \\
\bottomrule
\end{tabular}}
\end{center}

\section*{Osmosis}

\subsection*{Water as a solvent}

Water is a very good \textbf{solvent} 溶剂 — many substances \textbf{dissolve} 溶解 in it to make a \textbf{solution} 溶液. This matters for \textbf{digestion} 消化 (food must dissolve before it can be used), \textbf{excretion} 排泄 (wastes are carried away dissolved in water) and \textbf{transport} 运输 (substances travel around the body dissolved in blood and sap).

\subsection*{What osmosis is}

\textbf{Osmosis} 渗透 is a special kind of diffusion: the net movement of water across a \textbf{partially permeable membrane} 半透膜. This kind of membrane has tiny holes that let small water molecules through but hold back larger solute particles. Water moves into and out of cells by osmosis through the cell membrane.

A \textbf{concentrated} solution has a lot of dissolved solute and little water. A \textbf{dilute} solution has little solute and a lot of water.

\textbf{(Supplement)} Osmosis is the net movement of water molecules from a region of \textbf{higher water potential} 水势 (a dilute solution — more water) to a region of lower water potential (a concentrated solution — less water), through a partially permeable membrane.

\subsection*{Osmosis and plant cells}

When you put plant tissue into different solutions, water moves by osmosis:

\begin{itemize}
\item In \textbf{pure water or a dilute solution}, water moves \textbf{into} the cells. The cells swell and become firm, or \textbf{turgid} 膨胀. The water presses outward on the cell wall; this outward push is the \textbf{turgor pressure} 膨压. Turgid cells make a plant stand up straight — this is how plants are supported.

\item In a \textbf{concentrated solution}, water moves \textbf{out} of the cells. The cells lose their firmness and become soft, or \textbf{flaccid} 松软.

\item \textbf{(Supplement)} If even more water leaves, the cell membrane pulls away from the cell wall. This is \textbf{plasmolysis} 质壁分离.

\end{itemize}

You can investigate osmosis using \textbf{dialysis tubing} 透析管 (an artificial partially permeable membrane), or using cylinders cut from a potato. Measure their length or mass before and after soaking. Cylinders in pure water or a dilute solution gain length and mass; cylinders in a concentrated solution lose length and mass; in a solution of equal concentration there is no change.

\section*{Active transport}

\textbf{Active transport} 主动运输 is the movement of particles through a cell membrane from a region of \textbf{lower} concentration to a region of \textbf{higher} concentration — that is, \textbf{against} the concentration gradient. Because this is "uphill", it needs energy from \textbf{respiration} 呼吸作用.

\textbf{(Supplement)} Active transport lets a cell take in useful molecules or ions even when they are already more concentrated inside the cell. For example, \textbf{root hairs} 根毛 take up mineral ions from the soil by active transport. \textbf{Protein carriers} 载体蛋白 in the membrane pick up the molecules or ions and carry them across, using energy.

\section*{Comparing the three processes}

\begin{center}
\adjustbox{max width=\linewidth}{%
\begin{tabular}{llll}
\toprule
\textbf{Process} & \textbf{Direction} & \textbf{Energy from respiration?} & \textbf{What moves} \\
\midrule
diffusion & high → low concentration & no & particles, gases, solutes \\
osmosis & high → low water potential & no & water only \\
active transport & low → high concentration (against the gradient) & yes & molecules and ions \\
\bottomrule
\end{tabular}}
\end{center}

\section*{Exam tips}

\begin{itemize}
\item Always write \textbf{net} movement: particles move both ways, but the overall flow is down the gradient (for diffusion and osmosis).

\item Osmosis moves \textbf{water only}, and \textbf{only} through a partially permeable membrane. Diffusion can move many kinds of particle.

\item Only \textbf{active transport} uses energy from respiration, and only it goes \textbf{against} the gradient.

\item For the potato or plant experiment: water in → turgid, longer and heavier; water out → flaccid, shorter and lighter; equal concentration → no change.

\item A larger surface area, higher temperature, steeper concentration gradient and shorter distance all make diffusion faster.

\end{itemize}


\end{document}
